\section{The pricing test}
\label{sec:pricing-test}
\label{sec:pricing-protocol}

The pricing test splits an objective's effect on a metric into the value of
its change in a statistic and the cost of all else it alters, by imposing that
change on the baseline's samples (Fig.~\ref{fig:teaser}(a)).

\subsection{Definition}
\label{sec:pricing-def}

Let $B$ and $A$ be sets of $n$ samples from a baseline model and from a model
trained with the objective under test, drawn with the same sampler and
sampling seed, and let $M(X)$ be a metric that scores a set $X$ against a
reference set $R$, lower being better (FID throughout). The objective claims to
change a statistic $s(X) \in \mathbb{R}^K$ of a set of images and achieves
$\delta_A = s(A) - s(B)$. An operator $T$ (not the number of timesteps of
Section~\ref{sec:background}) acts on each image of a set, and $T_\delta$
denotes an operator constructed so that $s(T_\delta(X)) = s(X) + \delta$. The
value of an operator at the operating point $X$, the mechanism residual gap
(MRG) and the accounting that links them are
\begin{linenomath}
\begin{align}
  V(T; X) &= M(X) - M\big(T(X)\big), \label{eq:value}\\
  \mathrm{MRG} &= M(A) - M\big(T_{\delta_A}(B)\big), \label{eq:mrg}\\
  D = M(A) - M(B) &= \mathrm{MRG} - V, \qquad V = V(T_{\delta_A}; B),
  \label{eq:accounting}
\end{align}
\end{linenomath}
where $D$ is the metric gap and $V$, the value (or \emph{price}) of the
objective's change, is what imposing that change on the baseline's samples is
worth in $M$ ($V > 0$ means that the change lowers $M$, and $D > 0$ that $A$
is worse than $B$); the \emph{price curve} is the value as a function of the change
an operator produces. Equation~\eqref{eq:accounting} is an identity for any
operator; the preconditions of Section~\ref{sec:pricing-pre} give its terms
their meaning. The MRG, $D + V$, is the cost of everything the objective
alters besides $s$, per-image, anisotropic and phase changes included, so it
measures how much of the effect the stated mechanism fails to explain, not
what explains it. A positive MRG means that the objective is worse than the
baseline carrying the same change, not that this cost exceeds what the change
earns (that is $D > 0$); a negative MRG means that the objective gains more
than its change explains. The value is measured at a trained model's
operating point, not on real images, where the metric's first-order response
vanishes in the population limit.

\subsection{Statistic, operators and scoring}
\label{sec:pricing-ops}

\paragraph{Statistic.}
With $\hat y$ the centred 2-D DFT of the luminance $y$ of an image $x$ and
$U_k$ the frequencies $u$ at integer radius $r(u) = k$ from the zero
frequency, the statistic is the geometric-mean radial spectrum
\begin{equation}
  s_k(X) = \frac{1}{|X|} \sum_{x \in X} \frac{1}{|U_k|}
           \sum_{u \in U_k} 20 \log_{10} \lvert \hat y(u) \rvert
  \quad \text{(dB)},
  \label{eq:statistic}
\end{equation}
computed on all $n$ samples and all \num{\nRefImages} reference images
(Appendix~\ref{app:scoring}). The bands are the three index thirds of the
bins, on $32 \times 32$ images radii $0$--\bandLowMax{} (low),
\bandMidMin--\bandMidMax{} (mid) and \bandHighMin--\bandHighMax{} (high), and
a band value $\bar s_\beta$ is the unweighted mean of its bins. The axial
Nyquist radius is $16$, so the high band lies mostly in the diagonal corners
of the spectrum; its four outermost bins, with \nCoefCorner{} of its
\nCoefHigh{} coefficients, carry half the weight of the band value, a choice
on which $V$ and the MRG do not depend. The high-band deficit of $X$ is
$\bar s_{\mathrm{high}}(R) - \bar s_{\mathrm{high}}(X)$.

\paragraph{Operators.}
A spectral operator multiplies the centred DFT of each colour channel by one
real, positive, radially symmetric gain $g$ for all images and channels of a
set, $T_g(x)_c = \operatorname{Re} F^{-1}(g \odot F x_c)$, and leaves phases,
orientation and the spread of profiles across images untouched. Since the
luminance is linear in the channels, the gain
$g_\delta(u) = 10^{\delta_{r(u)}/20}$ moves $s$ by exactly $\delta$ before
clipping; it imposes the luminance change on every channel, and in bin~0 it
rescales mean brightness. The \emph{transplant} $T_{\delta_i}$, with
$\delta_i = s(A_i) - s(B_i)$ for the runs of seed $i$, is the canonical
operator because it reproduces the achieved change in every bin, not because
it is its only realisation; the \emph{reverse transplant} $T_{-\delta_i}$
removes the same change from $A_i$. The \emph{boost}, a raised-cosine
high-frequency gain of dose $d$~dB, is the instrument for the price curve;
its ramp begins at two thirds of the axial Nyquist frequency, inside the mid
band, so it also raises the upper mid-band bins slightly. The \emph{match}
$T_\delta$, $\delta = s(R) - s(B_i)$, removes the mean radial deficit in
every bin, a scale reference rather than an upper bound
(Section~\ref{sec:second-operator}). \emph{Band-restricted} transplants and
matches set $\delta$ to zero outside chosen bands, so the value of one band
is its effect with the others at the baseline's level. \emph{White noise},
uniform within $\pm 0.5$ or $\pm 1$ grey level per pixel and channel (at
$\pm 0.5$ the distribution of 8-bit rounding error), is the control for P0.
Outputs are clipped to $[0, 255]$, not rounded, and scored in floating point
with the resizer and feature extractor of \texttt{clean-fid}
\citep{parmar2022aliased} (Appendix~\ref{app:scoring}).

\subsection{Estimation and decision rule}
\label{sec:pricing-est}

Seed $i = 1, \dots, I$ fixes the initialisation and data order of both arms,
and all runs share one sampling seed, so $B_i$ and $A_i$ differ by the
objective alone (the two arms trained with an exponential moving average of
the weights, EMA, priced against MSE without it, also by weight averaging;
Section~\ref{sec:worth-arms}). The per-seed $D_i$, $V_i$ and
$\mathrm{MRG}_i$, with $D_i = \mathrm{MRG}_i - V_i$ exactly, are estimated by
their means with 95\% $t$-intervals and tested by two-sided paired
$t$-tests; since $\delta_i$ comes from the pair it is applied to, the
intervals include the uncertainty in the achieved change. When several
objectives are priced, the verdict rests on the Holm-adjusted $p$-values for
$\mathrm{MRG} = 0$ within the declared family (Section~\ref{sec:protocol}),
reported with the raw ones. The test of $\mathrm{MRG} = 0$ compares $A_i$
with $T_{\delta_i}(B_i)$, without $M(B_i)$, and is more precise than a test of
$D$ against the margin $-\bar V$ when
$\operatorname{Cov}(D_i, V_i) < -\mathrm{Var}(V_i)/2$, as $M(B_i)$ enters
$D_i$ and $V_i$ with opposite signs. An interval that contains zero bounds the
MRG and does not show that it is zero.

\begin{algorithm}[tbp]
\caption{The pricing test (same $n$ and $R$, unquantised scoring).}
\label{alg:pricing}
\begin{algorithmic}[1]
\Require seed pairs $(B_i, A_i)_{i=1}^{I}$; reference $R$; statistic $s$;
  metric $M$; declared family
\Ensure intervals and $p$-values of $\bar V$, $\overline{\mathrm{MRG}}$ and
  $\bar V_A - \bar V$; verdict; $V^\ast$
\For{$i = 1, \dots, I$}
  \State $\delta_i \gets s(A_i) - s(B_i)$ on all $n$ samples
    \Comment{P1}
  \State $B_i' \gets T_{\delta_i}(B_i)$, $A_i' \gets T_{-\delta_i}(A_i)$,
    unrounded; check $s(B_i') \approx s(A_i)$ per bin
    \Comment{P0, P2}
  \State $V_i \gets M(B_i) - M(B_i')$; \
    $\mathrm{MRG}_i \gets M(A_i) - M(B_i')$; \
    $V_{A,i} \gets M(A_i') - M(A_i)$
\EndFor
\State paired $t$-tests of $V_i$, $\mathrm{MRG}_i$ and $V_{A,i} - V_i$
  \Comment{P5}
\State price curve of the boost on $B$; white noise; other operators
  \Comment{P4, P0, P3}
\State verdict at $p < 0.05$, Holm-adjusted if several objectives: MRG
  $> 0$, $A$ falls short of the value of its change; MRG $< 0$, $A$ gains more
  than its change is worth; otherwise unresolved
\State tipping point of the identifying assumption:
  $V^\ast \gets -D_{\mathrm{lo}}$
\end{algorithmic}
\end{algorithm}

\subsection{Preconditions and controls}
\label{sec:pricing-pre}

Six preconditions give $V$ and the MRG their meaning;
Table~\ref{tab:preconditions} lists where each is checked and whether it
holds. As $\mathrm{MRG}_i = D_i + V_i$ with $D_i$ fixed by the arms, an upward
bias in $V$ produces a false positive, a shortfall the objective does not
have, and a downward bias a false negative.
\begin{description}
  \setlength{\itemsep}{1pt}\setlength{\parskip}{0pt}
  \item[P0] \emph{No re-quantisation.} Rounding to 8 bits adds a nearly white
    error field, strongest relative to image spectra in the high band. Its
    effect on the metric, hazard H1 (Section~\ref{sec:worth-curve}), biases
    $V$ upwards, so operators are scored in floating point; a gain still
    rescales $B$'s rounding noise.
  \item[P1] \emph{The objective changes the statistic it names, under the
    evaluation protocol}: $s$ is fixed before seeing what moved, and
    $\delta_A$ is measured on the samples that $M$ scores. For AWWL this
    holds under 100-step DDIM sampling \citep{song2021denoising} but not
    under 1000-step ancestral sampling (Section~\ref{sec:mechanism}).
  \item[P2] \emph{$\delta$ is measured on $A$, applied to $B$ and measured
    back} in every bin of the luminance statistic: the transplant satisfies
    this up to clipping, the boost only in the high-band mean.
  \item[P3] \emph{More than one operator.} Operators that impose the same
    $\delta$ can differ in what else they change (phase coherence, per-image
    spread, chroma), so $V$ is specific to the operator. All spectral
    operators here are coherent: each multiplies every image's spectrum by
    the same deterministic radial gain and so keeps its phases, and only the
    transplant imposes $\delta_A$. White noise adds incoherent energy
    (random phases, independent across pixels and images), which the metric
    can value far more per decibel of change in $s$
    (Section~\ref{sec:worth-checks}); a mismatch with training can therefore
    bias $V$ either way, and side effects such as ringing bias it downwards.
  \item[P4] \emph{The metric registers $s$}: the price curve is not flat at
    the operating point; otherwise $V \approx 0$ cannot separate a worthless
    change from an invisible one (Section~\ref{sec:dreambooth}).
  \item[P5] \emph{No interaction}: removing the change from $A$'s samples
    costs what adding it to $B$'s earns,
    $V_A = M(T_{-\delta_A}(A)) - M(A) \approx V$, judged by the paired
    difference. Otherwise Eq.~\eqref{eq:accounting} admits more than one
    split, and a change worth less at $A$ than at $B$ biases $V$ upwards.
\end{description}

Two controls run through Algorithm~\ref{alg:pricing} unchanged. The positive
control prices the gain of longer training, which targets nothing spectral,
and must return a negative MRG; since $\mathrm{MRG} = D + V$, it fails only
if the transplant is worth at least the whole training gain, so it excludes
only a gross upward bias in $V$. The null control prices a static band
weighting that leaves $s$ essentially unchanged, so $V$ must be close to zero
and the MRG equal to the arm's FID gap. When $|V|$ is small, as in both
controls, $\mathrm{MRG} \approx D$, so they check the verdict's direction and
the accounting, not the identifying assumption that $M$ values an
operator-made change as it would the same change made by training. A
planted-mechanism control tests that assumption directly: a baseline trained
on data that carry $\delta_A$ acquires the change through training alone and
must return an MRG of zero. Where no such arm is trained, a sensitivity
analysis reports how far the assumption may fail. If the training-made change is worth a constant
$V_{\mathrm{train}}$, the cost outside the change, $D + V_{\mathrm{train}}$,
keeps its uncorrected 95\% interval above zero whenever $V_{\mathrm{train}}$ exceeds the
tipping point $V^\ast = -D_{\mathrm{lo}}$ ($D_{\mathrm{lo}}$ the lower bound
of $D$'s interval), i.e.\ exceeds the fraction $V^\ast / V$ of the
operator's value (Section~\ref{sec:worth-checks}).

\subsection{Mediation reading}
\label{sec:pricing-reading}

In mediation analysis \citep{robins1992identifiability,pearl2001direct}, the
objective is the treatment, $s$ the mediator and $M$ the outcome, with the
training seed as the unit and $s$ and $M$ functionals of its sample set. With
$M(z, v)$ the metric of samples from objective $z \in \{0, 1\}$ (baseline,
objective under test) with statistic $v$, and $s_z$ the statistic that $z$
produces, $M(T_{\delta_A}(B))$ estimates $M(0, s_1)$ and
$M(T_{-\delta_A}(A))$ estimates $M(1, s_0)$. The MRG is then analogous to the
total direct effect $M(1, s_1) - M(0, s_1)$
\citep{robins1992identifiability}, and $-V$ to the pure indirect effect
$M(0, s_1) - M(0, s_0)$. The analogy is not an equivalence, since an operator
counterfactual is not a potential outcome of a trained model. It rests on
three assumptions: (i)~the operator sets the mediator of $B$ to $s_1$ and
changes nothing else that $M$ registers (P0, P2); (ii)~the identifying
assumption, untested but addressed by the sensitivity analysis; (iii)~no
treatment--mediator interaction, under which the total direct effect equals
the natural direct effect $M(1, s_0) - M(0, s_0)$ \citep{pearl2001direct};
their difference, $V - V_A$, is what P5 measures. As the mediator is set by
intervention, (i) and (ii) replace the usual assumption of no unmeasured
mediator--outcome confounding.
